\section{Cumulants versus moments}
\label{app:cumulants}

Joint cumulants are generated by the logarithm of the moment generating function. By expanding these generating functions, one arrives at the moment-cumulants formula~\cite{leonov1959}, which specifies how a mixed moment is expanded in terms of joint cumulants,
\begin{align}
  \EE[X_{1}\cdots X_{n}]=\sum_{p\in P_{n}}\prod_{b\in p}\,\EE^{c}\!\!\left[ \prod_{i\in b} X_{i}\right]\,,\label{eq:5}
\end{align}
where $P_{n}$ is the set of partitions of $\{1,\dots,n\}$, hence $b$ is a block in the partition $p$. For instance,
\begin{align}
  \EE[X_{1}X_{2}X_{3}]&=\EE^{c}[X_{1},X_{2},X_{3}]+\EE^{c}[X_{1},X_{2}]\EE^{c}[X_{3}]+\EE^{c}[X_{1},X_{3}]\EE^{c}[X_{2}]+\EE^{c}[X_{2},X_{3}]\EE^{c}[X_{1}]\nonumber\\
  &\qquad+\EE^{c}[X_{1}]\EE^{c}[X_{2}]\EE^{c}[X_{3}]\,.
\end{align}
We will use the decomposition~\eqref{eq:5} to expand mixed moments in terms of cumulants which we then compute using Feynman diagrams.

Intuitively, joint cumulants correspond to subtracting all factorizations of the mixed moment, e.g.
\begin{align}
  \EE^{c}[X_{1},X_{2},X_{3}]&=\EE[X_{1}X_{2}X_{3}]-\EE[X_{1}X_{2}]\EE[X_{3}]-\EE[X_{1}X_{3}]\EE[X_{2}]-\EE[X_{2}X_{3}]\EE[X_{1}]\nonumber\\
  &\qquad+2\EE[X_{1}]\EE[X_{2}]\EE[X_{3}]\,,
\end{align}
where the last term is added twice to compensate for the factorizations in the three terms before. In general, a joint cumulant is given in terms of mixed moments by~\cite{leonov1959}
\begin{align}
  \EE^{c}[X_{1},\cdots, X_{n}]=\sum_{p\in P_{n}}(|p|-1)!(-1)^{|p|-1}\prod_{b\in p}\EE\left[ \prod_{i\in b}X_{i} \right]\,,
\end{align}
where $|p|$ is the number of blocks in the partition $p$.


%%% Local Variables:
%%% mode: LaTeX
%%% TeX-master: "ntk_feynman"
%%% End:
